\documentclass[12pt]{article}
\usepackage[T1]{fontenc}
\usepackage[utf8]{inputenc}
\usepackage{lmodern}
\usepackage{textcomp}
\usepackage{enumitem}
\usepackage{geometry}
\geometry{margin=1in}
\begin{document}
\begin{center}
Answer Key - Geography - K-12 Level Assessment 21 \
\small (Answers are ordered exactly as in the question paper.)
\end{center}

\section*{Multiple Choice}
\begin{enumerate}[label=\arabic*.]
\item \textbf{Answer:} B
\\ \textit{Options:} A) A high level of confidence in central government; B) The existence of strong separatist groups; C) The existence of national transportation networks; D) The national anthem
\\ \textit{Note:} The existence of strong separatist groups represents a force that can lead to the fragmentation of a state, rather than a centripetal force that promotes unity and cohesion within a state.
\item \textbf{Answer:} A
\\ \textit{Options:} A) high birth rates with high but fluctuating death rates.; B) declining birth rates with continuing high death rates.; C) low birth rates with continuing high death rates.; D) high birth rates with declining death rates.
\\ \textit{Note:} The first stage of the demographic transition is characterized by high birth rates and high but fluctuating death rates due to limited access to healthcare, poor sanitation, and subsistence farming, leading to a stable population size with minimal growth.
\item \textbf{Answer:} D
\\ \textit{Options:} A) a water divider.; B) a watercourse.; C) an artificial boundary.; D) a natural boundary.
\\ \textit{Note:} The Rio Grande is considered a natural boundary because it is a physical geographic feature formed by the river itself that serves as a dividing line between the two countries.
\item \textbf{Answer:} C
\\ \textit{Options:} A) Calcutta.; B) Bombay.; C) Cambodia.; D) Bali.
\\ \textit{Note:} The largest Hindu temple complex, Angkor Wat, is located in Cambodia and is renowned for its vast scale and intricate architecture, making it the most significant and expansive religious site dedicated to Hinduism.
\item \textbf{Answer:} B
\\ \textit{Options:} A) alpine tundra.; B) tropical forests.; C) flood plains.; D) deserts.
\\ \textit{Note:} Shifting cultivation is most often practiced in tropical forests because these regions have nutrient-rich soils and a warm, humid climate that supports the growth of diverse vegetation, making them suitable for slash-and-burn agricultural techniques.
\end{enumerate}

\section*{True / False}
\begin{enumerate}[label=\arabic*.]
\setcounter{enumi}{5}
\item \textbf{Answer:} True
\\ \textit{Options:} A) True; B) False
\\ \textit{Note:} Japan is the least self-sufficient in natural resources because it relies heavily on imports for essential materials such as oil, minerals, and timber to sustain its industrial economy.
\item \textbf{Answer:} True
\\ \textit{Options:} A) True; B) False
\\ \textit{Note:} The statement is true because the influx of new residents often brings different cultural backgrounds, economic status, and social networks, which can alter the existing community dynamics and relationships among residents.
\item \textbf{Answer:} True
\\ \textit{Options:} A) True; B) False
\\ \textit{Note:} The Industrial Revolution introduced new technologies and machinery that improved agricultural productivity and efficiency, thereby facilitating the advancements of the Second Agricultural Revolution.
\item \textbf{Answer:} False
\\ \textit{Options:} A) True; B) False
\\ \textit{Note:} A general purpose map displays a variety of information, including physical features, political boundaries, and demographic data, rather than focusing exclusively on one type of statistic.
\item \textbf{Answer:} False
\\ \textit{Options:} A) True; B) False
\\ \textit{Note:} East Asia's urbanization, while significant, is not faster than that of regions like sub-Saharan Africa and parts of South Asia, which are currently experiencing even more rapid urban growth due to high population growth rates and migration trends.
\end{enumerate}

\section*{Long Form Response}
\begin{enumerate}[label=\arabic*.]
\setcounter{enumi}{10}
\item \textbf{Answer:} A country's agriculture is influenced by several environmental factors, including climate, soil quality, water availability, and topography. Climate affects which crops can grow in a region, while soil type determines the nutrients available for plant growth. Water availability is crucial for irrigation, especially in arid areas, and topography can impact farming practices and crop yields. While the amount of fertilizer produced in a country can support agricultural productivity, it is not considered a crucial factor because effective farming relies more on the natural environment and sustainable practices. If the soil and climate are unsuitable, even excessive fertilizer use cannot compensate for poor growing conditions.
\\ \textit{Note:} correct because it identifies key environmental factors such as climate, soil quality, water availability, and topography that directly influence agricultural practices and yields, while clarifying that fertilizer production, though supportive, is not as critical as these fundamental factors.
\item \textbf{Answer:} The relationships and associations between various North American regions and their demographic groups reveal a complex picture that often contradicts common perceptions. For instance, while many people associate African Americans primarily with the Southern United States, there is a significant presence of African American communities in the Southern prairie provinces of Canada, such as Alberta and Saskatchewan. This demographic group has historical roots in these areas, often linked to the migration of black settlers in the late 19 th and early 20 th centuries. Additionally, the misconception that African Americans are only concentrated in urban areas overlooks their contributions and presence in rural settings across North America. Understanding these dynamics helps to appreciate the rich tapestry of cultural and demographic diversity in the continent.
\\ \textit{Note:} correct because it illustrates the nuanced and often overlooked presence of African American communities outside their commonly perceived geographic boundaries, emphasizing the importance of recognizing demographic diversity across North America.
\end{enumerate}

\vfill
\noindent\textit{End of Answer Key}
\end{document}